\chapter{The Rule~4 Audit: Are the Excavator False Alarms Real?}
\label{chap:s2-7}

\section{Goal}

The \S5.5 audit (Chapter~\ref{chap:s2-6}) found that rule~4 is the rule the annotation most often leaves
unmarked: a worker within reach of an operating excavator. Every rule~4 accuracy figure in Study~2 is
scored against that incomplete reference. Two readings depend on it. The unquantised model's rule~4
``false alarms'' (31\% of compliant images) may be real hazards. Its much higher rule~4 sensitivity (81\%
against 26\% for 4-bit) is measured against the same incomplete reference. This chapter relabels a
stratified draw of the sample's rule~4 questions and corrects both tiers' rule~4 confusion matrix.

\section{What Was Done}

\textbf{The draw (decided 2026-11-12, committed as \texttt{1973153} before any image was viewed).} The
frozen sample holds 323 rule~4 questions, one per image. \texttt{15\_rule4\_audit\_sample.py} splits them
into nine cells by stratum and by what each tier answered:
\begin{itemize}
\item C: compliant images;
\item H: hazard images where rule~4 is not annotated;
\item P: images where rule~4 is annotated.
\end{itemize}
Each cell is drawn with its own seeded generator (12 per negative cell, 8 per positive cell; small cells
whole), giving 92 images. The positive cells are in the draw so the reader cannot infer the answer from
the draw itself, and they estimate how often a marked rule~4 violation is wrong.

\textbf{The reader.} As in Chapter~\ref{chap:s2-6}, Claude (Opus~5.5) answered first. It answered rule~4
only, in a seeded random order, against the dataset's rule text, without seeing the cell, the annotation,
or either model's answer. The answers were committed (\texttt{687059b}) before scoring. The standard
applied the owner's two corrections from the \S5.5 review:
\begin{itemize}
\item camera angle shortens distances;
\item the excavator must be working or visibly manned.
\end{itemize}
Where either could not be judged, the answer is ``can't tell''.

\textbf{The correction.} \texttt{16\_rule4\_audit\_score.py} takes, per cell, the share of answered rows the
reader calls a violation and gives every one of the 323 rows its cell's share as the probability that it
is a real violation. Each tier's confusion matrix is then recomputed in expectation, beside the
annotation-based one, with a bootstrap over answered rows within each cell.

\section{The Owner's Review}

The owner reviewed every image where the first reader's answer disagreed with the annotation or was
``can't tell'': 40 of the 92. The other 52 agree with the annotation and keep the first reader's answer.
\texttt{18\_rule4\_adjudicate.py} merges the review, with the owner's answer taking precedence. The owner
changed 8:
\begin{itemize}
\item \emph{First reader Yes, owner No or can't tell (5 of 25).} The worker is outside the radius
(0007497, 0010299) or probably outside it (0004350). The excavator may not be working (0010103). One
note asks whether a mini excavator is a threat at all (0012295).
\item \emph{First reader can't tell, owner decided (3 of 15).} Two No (0001436, 0003886), one Yes
(0002534: within four or five metres of the hammer).
\end{itemize}
The owner kept 20 of the first reader's 25 unannotated Yes calls. Twelve images stay ``can't tell'' for
both readers. On one (0005496) the operator's cab is out of frame, so whether it is manned cannot be
judged. On another (0010267) the workers stand above the lowered bucket, a case the rule's text does not
cover.

\section{Results (adjudicated)}

\begin{table}[tbp]
\centering
\begin{tabular}{lrrrr}
\toprule
cell & size & drawn & says violation & first reader \\
\midrule
C, both tiers No & 85 & 12 & 1 / 12 & 1 / 12 \\
C, 4-bit No, bf16 Yes & 32 & 12 & 5 / 10 & 7 / 10 \\
C, 4-bit Yes & 8 & 8 & 3 / 7 & 3 / 7 \\
H, both No & 28 & 12 & 2 / 9 & 2 / 9 \\
H, 4-bit No, bf16 Yes & 59 & 12 & 6 / 12 & 5 / 10 \\
H, both Yes & 41 & 12 & 4 / 9 & 7 / 11 \\
P (annotated), all cells & 70 & 24 & 18 / 18 & 18 / 18 \\
\bottomrule
\end{tabular}
\caption{Rule~4 audit by draw cell, adjudicated, beside the first reader alone. ``Can't tell'' (15 of 92
in both) is excluded from the share.}
\label{tab:s27cells}
\end{table}

\begin{table}[tbp]
\centering
\begin{tabular}{lrrrr}
\toprule
 & \multicolumn{2}{c}{4-bit} & \multicolumn{2}{c}{bf16} \\
 & annotation & corrected & annotation & corrected \\
\midrule
sensitivity & 25.7\% & 26.4\% [18.3, 34.5] & 81.4\% & 82.2\% [72.8, 90.0] \\
specificity & 80.6\% & 84.1\% [77.1, 91.4] & 45.1\% & 58.1\% [49.8, 68.2] \\
false alarm, compliant images & 6.4\% & 4.6\% [1.3, 7.6] & 31.2\% & 20.3\% [10.8, 28.6] \\
false alarm, other-hazard images & 32.0\% & 30.8\% [15.1, 45.7] & 78.1\% & 70.6\% [56.9, 81.1] \\
\bottomrule
\end{tabular}
\caption{Rule~4 over the 323 sample rows, scored against the annotation and corrected by the adjudicated
relabels (\texttt{results/rule4\_audit\_score\_owner-adjudicated.txt}). The first reader alone had put
bf16's compliant-image false-alarm rate at 14.8\% and its specificity at 63.3\%.}
\label{tab:s27}
\end{table}

\textbf{The annotation misses more rule~4 violations than it marks.} Every annotated rule~4 violation
that could be judged was confirmed (18 of 18). The adjudicated labels put about 80 real violations among
the 253 unannotated rows, against 70 annotated. The first reader alone had estimated 95. Most of them
are on images where the unquantised model said Yes: 5 of 10 such compliant images and 6 of 12 such
other-hazard images.

\textbf{What changes for the two tiers.}
\begin{itemize}
\item \emph{bf16.} About a third of its ``false alarms'' on compliant images are real hazards. Its
corrected false-alarm rate there is 20.3\%, not 31.2\%, and its rule~4 specificity rises from 45\% to
58\%.
\item \emph{4-bit.} Its rule~4 blindness is real. Corrected sensitivity is 26\%, and the gap to bf16
(82\%) is unchanged.
\end{itemize}
So quantisation costs rule~4 sensitivity, as Chapter~\ref{chap:s2-5} found, and a substantial part of
what looked like the unquantised model's over-triggering is really the annotation's omission. The owner's
review moved the figures back toward the annotation without changing either conclusion.

\textbf{Test--retest.} Four images were also in the \S5.5 draw. The reader gave the same rule~4 answer on
three. On the fourth (0009901) it moved from Yes to ``can't tell'', applying the stricter standard
adopted after the owner's review.

\section{Standard vs.\ Non-Standard Choices}

\textbf{Standard.} Stratified relabelling with design weights. Strata fixed and committed before any
image was viewed.

\textbf{Non-standard, and stated in the paper.}
\begin{itemize}
\item The cells are defined by the models' answers, which is what lets the correction apply to each tier.
The reader was kept blind to the cells for exactly this reason.
\item The correction assumes that the share of real violations is the same for every row in a cell. Cells
of 8 to 12 answered rows make that share coarse, so the intervals are wide.
\item The first reader is a vision-language model. The owner reviewed only the 40 images where it
disagreed with the annotation or could not tell; the 52 where it agreed were not re-read. The owner kept
20 of its 25 unannotated Yes calls, close to the \S5.5 review, where 9 of 11 were kept.
\item Fifteen images (16\%) stay ``can't tell''. They are left out of each cell's share, which amounts to
assuming they are violations at the same rate as the answered rows.
\end{itemize}

\section{Implications}

The rule~4 row of every Study~2 table reads differently once the reference is corrected. The 4-bit
model's near-blindness to rule~4 (F2) stands. The unquantised model's rule~4 false-alarm rate on
compliant images is about two thirds of what the annotation says (20\% against 31\%). More generally, a model that is more sensitive
than the annotators is penalised as a false-alarm generator, which matters for how the paper reads F1
(``accuracy is bought with false alarms'').
